\documentclass[10pt,a4paper]{amsart}
\usepackage{amsmath,amssymb,amsthm,graphicx}
\usepackage{verbatim}
\newenvironment{reference}{\comment}{\endcomment}
%\newenvironment{reference}{}{}
\newenvironment{slogan}{\comment}{\endcomment}
\newenvironment{history}{\comment}{\endcomment}

% For commutative diagrams we use Xy-pic
\usepackage[all]{xy}

% We use 2cell for 2-commutative diagrams.
\xyoption{2cell}
% Theorem environments.
%
\theoremstyle{plain}
\newtheorem{theorem}[subsection]{Theorem}
\newtheorem{proposition}[subsection]{Proposition}
\newtheorem{lemma}[subsection]{Lemma}

\theoremstyle{definition}
\newtheorem{definition}[subsection]{Definition}
\newtheorem{example}[subsection]{Example}
\newtheorem{exercise}[subsection]{Exercise}
\newtheorem{situation}[subsection]{Situation}

\theoremstyle{remark}
\newtheorem{remark}[subsection]{Remark}
\newtheorem{remarks}[subsection]{Remarks}

\numberwithin{equation}{subsection}

% Macros
%
\def\lim{\mathop{\mathrm{lim}}\nolimits}
\def\colim{\mathop{\mathrm{colim}}\nolimits}
\def\Spec{\mathop{\mathrm{Spec}}}
\def\Hom{\mathop{\mathrm{Hom}}\nolimits}
\def\Ext{\mathop{\mathrm{Ext}}\nolimits}
\def\SheafHom{\mathop{\mathcal{H}\!\mathit{om}}\nolimits}
\def\SheafExt{\mathop{\mathcal{E}\!\mathit{xt}}\nolimits}
\def\Sch{\mathit{Sch}}
\def\Mor{\mathop{\mathrm{Mor}}\nolimits}
\def\Ob{\mathop{\mathrm{Ob}}\nolimits}
\def\Sh{\mathop{\mathit{Sh}}\nolimits}
\def\NL{\mathop{N\!L}\nolimits}
\def\CH{\mathop{\mathrm{CH}}\nolimits}
\def\proetale{{pro\text{-}\acute{e}tale}}
\def\etale{{\acute{e}tale}}
\def\QCoh{\mathit{QCoh}}
\def\Ker{\mathop{\mathrm{Ker}}}
\def\Im{\mathop{\mathrm{Im}}}
\def\Coker{\mathop{\mathrm{Coker}}}
\def\Coim{\mathop{\mathrm{Coim}}}

% Boxtimes
%
\DeclareMathSymbol{\boxtimes}{\mathbin}{AMSa}{"02}

%
% Macros for moduli stacks/spaces
%
\def\QCohstack{\mathcal{QC}\!\mathit{oh}}
\def\Cohstack{\mathcal{C}\!\mathit{oh}}
\def\Spacesstack{\mathcal{S}\!\mathit{paces}}
\def\Quotfunctor{\mathrm{Quot}}
\def\Hilbfunctor{\mathrm{Hilb}}
\def\Curvesstack{\mathcal{C}\!\mathit{urves}}
\def\Polarizedstack{\mathcal{P}\!\mathit{olarized}}
\def\Complexesstack{\mathcal{C}\!\mathit{omplexes}}
% \Pic is the operator that assigns to X its picard group, usage \Pic(X)
% \Picardstack_{X/B} denotes the Picard stack of X over B
% \Picardfunctor_{X/B} denotes the Picard functor of X over B
\def\Pic{\mathop{\mathrm{Pic}}\nolimits}
\def\Picardstack{\mathcal{P}\!\mathit{ic}}
\def\Picardfunctor{\mathrm{Pic}}
\def\Deformationcategory{\mathcal{D}\!\mathit{ef}}
\usepackage{cleveref}
\crefname{originallemma}{original lemma}{original lemmas}
\Crefname{originallemma}{Original lemma}{Original lemmas}
\makeatletter
\newlabel{fields-lemma-galois-goes-up}{{09I2}{source}{}{sourcepointer-1}{}}
\newlabel{fields-lemma-galois-goes-up@cref}{{[originallemma][0][]09I2}{[1][0][]source}{}{}{}}
\makeatother
\crefname{originallemma}{original lemma}{original lemmas}
\Crefname{originallemma}{Original lemma}{Original lemmas}
\makeatletter
\newlabel{fields-lemma-normal-closure-galois}{{0EXM}{source}{}{sourcepointer-2}{}}
\newlabel{fields-lemma-normal-closure-galois@cref}{{[originallemma][0][]0EXM}{[1][0][]source}{}{}{}}
\makeatother
\title{Injectivity of the degree zero cohomology along a finite separable extension}
\author{The Stacks Project Authors}
\date{stacks-project:weil.tex:local-snapshot-20261012; GNU FDL 1.2 or later}
\begin{document}
\maketitle
\noindent\textit{Editorial note.} All delivered mathematical text is the author's own, copied from the pinned Stacks Project chapter file weil.tex; no statement or proof step was rewritten, reordered or removed.\par\medskip
\noindent\textit{Editorial note.} The counted proof is the author's complete argument: the reduction to the Galois closure, the splitting of the product of the spectra through the tensor product decomposition, and the induced injectivity on cohomology.\par\medskip
\noindent\textit{Editorial note.} The label lemma-injective-H0 names the counted statement, and the proof environment immediately following it in the source is delivered verbatim as its proof.\par\medskip
\noindent\textit{Editorial note.} The counted statement is carried by the author's own lemma environment in the chapter source.\par\medskip
\noindent\textit{Editorial note.} Pre-open gate: python operations/agent-authoring/tools/precard\_gate.py --source weil.tex --label lemma-injective-H0 --start 5302 --end 5333 --local-id 28152 returned RESULT=OK (L28152 free, occupancy FREE, public interval clear). Pre-assembly tag self-check: canonical tag resolved and not present in reference/public-items.jsonl.\par\medskip
\noindent\textit{Editorial note.} The statement and proof are copied from the pinned chapter file; the cited inputs are kept as external pointers with their canonical Stacks tags.\par\medskip
\noindent\textit{Allowed standard prerequisites (editorial).}\begin{itemize}
\item Galois closure of a finite separable extension. Fields, the lemma that the normal closure of a finite separable extension is Galois, and the lemma that a Galois extension of a Galois extension remains Galois. Allows replacing the top field by its Galois closure over the base.
\item Splitting of the tensor product of a finite Galois extension. The decomposition of the tensor product of a finite Galois extension of a field with a field between the base and the extension into a product indexed by the Galois group. Produces the isomorphism between the disjoint union of copies of the spectrum of the top field and the product of the spectra, and hence the induced isomorphism on cohomology.
\end{itemize}\medskip
\ifdefined\hypertarget\hypertarget{sourcepointer-1}{}\fi
\noindent\textit{Source reference.} Original lemma 09I2: sources/stacks-chow/tags.txt, source lines 12025--12025. Prerequisite pointer invoked at proof line(s) 5314; it names a standard fact of a sibling chapter that the counted argument uses, and every step of the counted proof built on it is retained verbatim in the delivered body.\par\medskip
\ifdefined\hypertarget\hypertarget{sourcepointer-2}{}\fi
\noindent\textit{Source reference.} Original lemma 0EXM: sources/stacks-chow/tags.txt, source lines 18662--18662. Prerequisite pointer invoked at proof line(s) 5312; it names a standard fact of a sibling chapter that the counted argument uses, and every step of the counted proof built on it is retained verbatim in the delivered body.\par\medskip
\begin{lemma}
\label{lemma-injective-H0}
Assume given data (D0), (D1), and (D2') satisfying axioms (A1) -- (A7).
If $k''/k'/k$ are finite separable field extensions, then
$H^0(\Spec(k')) \to H^0(\Spec(k''))$ is injective.
\end{lemma}

\begin{proof}
We may replace $k''$ by its normal closure over $k$
which is Galois over $k$, see
Fields, Lemma \ref{fields-lemma-normal-closure-galois}.
Then $k''$ is Galois over $k'$ as well, see
Fields, Lemma \ref{fields-lemma-galois-goes-up}.
We deduce we have an isomorphism
$$
k' \otimes_k k'' \longrightarrow
\prod\nolimits_{\sigma \in \text{Gal}(k''/k')} k'',\quad
\eta \otimes \zeta \longmapsto (\eta \sigma(\zeta))_\sigma
$$
This produces an isomorphism
$\coprod_\sigma \Spec(k'') \to \Spec(k') \times \Spec(k'')$
which on cohomology produces the isomorphism
$$
H^*(\Spec(k')) \otimes_F H^*(\Spec(k''))
\to
\prod\nolimits_\sigma H^*(\Spec(k'')),\quad
a' \otimes a'' \longmapsto (\pi^*a' \cup \Spec(\sigma)^*a'')_\sigma
$$
where $\pi : \Spec(k'') \to \Spec(k')$ is the morphism
corresponding to the inclusion of $k'$ in $k''$.
We conclude the lemma is true by taking $a'' = 1$.
\end{proof}

\end{document}
